\chapter{Contexto histórico: dos epiciclos à gravitação}\label{cap:historia}

O resumo histórico a seguir é breve e segue as fontes usadas por \citeonline{farias2016}. Um tratamento mais extenso está em \cite{kuhn2002,martins2003,damasio2011,linton2004}.

\section{O problema dos planetas}

Para quem observa o céu a olho nu, as estrelas giram juntas, como se estivessem presas a uma esfera. Cinco pontos luminosos não seguem esse movimento. Mercúrio, Vênus, Marte, Júpiter e Saturno caminham devagar contra o fundo de estrelas, sempre perto da faixa do zodíaco, e de tempos em tempos param, recuam e voltam a avançar. Esse recuo é o \emph{movimento retrógrado}, e foi o grande desafio da astronomia antiga.

O próprio simulador reproduz o fenômeno. A \autoref{fig:marte-retrogrado} mostra a longitude eclíptica de Marte visto da Terra, calculada pelo programa a partir das condições iniciais de 1º de janeiro de 2026. Entre 10 de janeiro e 1º de abril de 2027 a longitude diminui: Marte anda para trás no céu. No plano de longitude e latitude, o caminho desenha o laço característico.

\begin{figure}[htb]
  \centering
  \caption{Movimento retrógrado de Marte em 2027, calculado pelo simulador}
  \label{fig:marte-retrogrado}
  \includegraphics[width=\textwidth]{marte_retrogrado.pdf}
  \legend{(a) Longitude eclíptica geocêntrica de Marte ao longo de dois anos; a faixa sombreada marca o intervalo em que ela decresce. (b) O mesmo intervalo no plano longitude--latitude, com o eixo de longitude invertido como no céu: o círculo e o quadrado marcam o início e o fim da retrogradação. Fonte: elaborada pelos autores com o integrador descrito no \autoref{cap:matematica}.}
\end{figure}

\begin{paracomecar}
O recuo é uma ilusão de perspectiva. A Terra anda mais depressa que Marte na órbita. Quando ela o ultrapassa, por algumas semanas Marte parece andar para trás, do mesmo jeito que um carro mais lento parece ir para trás quando o ultrapassamos na estrada.
\end{paracomecar}

\section{Ptolomeu: deferente, epiciclo, excêntrico e equante}

No \emph{Almagesto}, por volta de 150 d.C., Cláudio Ptolomeu reuniu a astronomia grega em um sistema matemático capaz de prever as posições dos planetas \cite{linton2004}. A Terra fica imóvel e cada planeta se move sobre um círculo pequeno, o \emph{epiciclo}, cujo centro percorre um círculo grande, o \emph{deferente}. A composição dos dois movimentos produz os laços retrógrados.

Dois refinamentos completam o modelo (\autoref{fig:ptolomeu}). O centro do deferente não coincide com a Terra: ele fica deslocado, num \emph{excêntrico}. E o centro do epiciclo não gira com velocidade angular constante em torno do centro do deferente nem em torno da Terra, mas em torno de um terceiro ponto, o \emph{equante}, simétrico à Terra em relação ao centro. O equante fazia o modelo concordar com as observações, mas violava o ideal de movimento circular uniforme, e foi isso que mais incomodou Copérnico \cite{farias2016,damasio2011}.

\begin{figure}[htb]
  \centering
  \caption{Elementos do modelo de Ptolomeu}
  \label{fig:ptolomeu}
  \begin{tikzpicture}[scale=1.15,>=Stealth]
    \def\R{2.4}
    \coordinate (C) at (0,0);
    \coordinate (T) at (-0.8,0);
    \coordinate (Q) at (0.8,0);
    \draw[thick,gray!70] (C) circle (\R);
    \node[gray!60!black,font=\small] at (-2.25,-1.95) {deferente};
    \coordinate (P) at (55:\R);
    \draw[thick,blue!60!black] (P) circle (0.75);
    \node[blue!60!black,font=\small,anchor=west] at ($(P)+(0.8,0.3)$) {epiciclo};
    \coordinate (Pl) at ($(P)+(150:0.75)$);
    \draw[blue!60!black] (P) -- (Pl);
    \fill[red!70!black] (Pl) circle (2pt) node[above left,font=\small] {planeta};
    \draw[dashed,orange!80!black] (Q) -- (P);
    \draw[gray] (C) -- (P);
    \fill (C) circle (1.5pt) node[above=2pt,font=\small] {centro};
    \fill[green!40!black] (T) circle (2.5pt) node[below left=1pt,font=\small] {Terra};
    \fill[orange!80!black] (Q) circle (1.8pt) node[below right=1pt,font=\small] {equante};
    \draw[<->,gray] ($(T)+(0,-0.75)$) -- node[below,font=\scriptsize]{excentricidade} ($(C)+(0,-0.75)$);
    \draw[->,orange!80!black] ($(Q)+(0.65,0)$) arc (0:60:0.65);
    \node[orange!80!black,font=\scriptsize,anchor=west] at ($(Q)+(0.7,0.25)$) {$\omega$ constante};
  \end{tikzpicture}
  \legend{O centro do epiciclo percorre o deferente com velocidade angular constante vista do equante, e não do centro nem da Terra. Fonte: elaborada pelos autores, a partir de \citeonline{linton2004} e \citeonline{farias2016}.}
\end{figure}

\section{Copérnico: o mesmo círculo, outro centro}

Nicolau Copérnico (1473--1543) apresentou o heliocentrismo primeiro num manuscrito curto, o \emph{Commentariolus}, que circulou por volta de 1514, e depois no \emph{De revolutionibus orbium coelestium}, publicado em 1543 \cite{martins2003,kuhn2002}. Com o Sol no centro, o movimento retrógrado deixa de exigir um epiciclo próprio para cada planeta: ele aparece naturalmente quando a Terra ultrapassa um planeta externo, ou é ultrapassada por um interno.

Copérnico, porém, não abandonou os círculos. Para reproduzir as variações de velocidade que Ptolomeu explicava com o equante, ele usou epiciclos menores. No \emph{Commentariolus}, cada planeta tem um deferente centrado no Sol e dois epiciclos. Os raios são dados em \emph{partes}, com o raio da órbita da Terra dividido em 25 partes \cite{martins2003,farias2016}. A \autoref{tab:copernico-raios} reproduz esses valores.

\begin{table}[htb]
  \centering
  \caption{Raios do deferente e dos epiciclos no \emph{Commentariolus}}
  \label{tab:copernico-raios}
  \begin{tabular}{lccc}
    \toprule
    Planeta & Deferente $R$ & 1º epiciclo $r_1$ & 2º epiciclo $r_2$ \\
    \midrule
    Mercúrio & $9 + 2/5$ & $1 + 41'$ & $34'$ \\
    Vênus    & $18$ & $3/4$ & $1/4$ \\
    Marte    & $38$ & $5 + 34'$ & $1 + 51'$ \\
    Júpiter  & $130 + 5/12$ & $10 + 6'$ & $3 + 22'$ \\
    Saturno  & $230 + 1/6$ & $19 + 41'$ & $6 + 34'$ \\
    \bottomrule
  \end{tabular}
  \legend{Valores em partes, com o raio da órbita terrestre igual a 25 partes; o símbolo $'$ indica sessenta avos de parte. O deferente de Marte segue a correção de Rosen adotada por \citeonline{farias2016}. Fonte: \citeonline{farias2016}, a partir de \citeonline{martins2003}.}
\end{table}

A tabela mostra dois fatos, explicados no \autoref{cap:resultados}. Primeiro, dividindo os deferentes por 25, obtêm-se as distâncias médias ao Sol em unidades astronômicas com erro de poucos por cento: Marte, $38/25 = 1{,}52$; Júpiter, $5{,}22$; Saturno, $9{,}21$. Segundo, em todos os planetas o primeiro epiciclo é \emph{exatamente três vezes} o segundo.

\section{Kepler e Newton}

Johannes Kepler (1571--1630), trabalhando com as observações de Tycho Brahe, abandonou os círculos e formulou as três leis que levam seu nome \cite{kuhn2002}:

\begin{enumerate}
  \item cada planeta descreve uma elipse com o Sol num dos focos;
  \item a reta que liga o Sol ao planeta varre áreas iguais em tempos iguais;
  \item o quadrado do período é proporcional ao cubo do semieixo maior.
\end{enumerate}

As duas primeiras foram publicadas na \emph{Astronomia nova}, em 1609, e a terceira na \emph{Harmonices mundi}, em 1619. Isaac Newton mostrou nos \emph{Principia}, em 1687, que as três decorrem de uma única lei de força, a gravitação universal, somada às leis do movimento. O \autoref{cap:fisica} refaz essa demonstração.

\section{Fourier e a pergunta do TCC}

Em 1807 Joseph Fourier apresentou o estudo da propagação do calor em que funções periódicas aparecem escritas como somas de senos e cossenos, publicado em 1822 na \emph{Théorie analytique de la chaleur} \cite{pifer2015}. Essa ideia dá uma resposta moderna para a eficácia dos epiciclos. Uma órbita fechada é uma curva periódica no plano. Se cada termo de uma série de Fourier for um círculo girando com velocidade constante, a série inteira é uma cadeia de epiciclos.

Essa é a pergunta de \citeonline{farias2016}: os epiciclos de Copérnico são uma aproximação da série de Fourier das órbitas reais? Para responder, a série precisa estar na forma complexa (\autoref{cap:matematica}). Escrita assim, ela reproduz a construção de Copérnico termo a termo, até primeira ordem na excentricidade (\autoref{sec:copernico-revisitado}).
